\chapter{The corpus}

\textbf{Purpose:} Record what each document in the corpus is, as established by opening it, and
separate that from what its filename says. \textbf{Scope:}
\texttt{Downloads/\allowbreak{}millenium/\allowbreak{}} as it stood on 2026-09-11, seven files,
7{,}852{,}267 bytes.

\section{What is in it}

Seven documents. Each was opened and identified from its own first page and embedded metadata, not
from its filename. Five are the official Clay Mathematics Institute problem descriptions. Two are
Riemann.

\begin{center}
\begin{tabular}{@{}llrr@{}}
\toprule
File & Identified as & Pages & Bytes \\
\midrule
\texttt{pvsnp.pdf}       & Stephen Cook, \emph{The P versus NP Problem}      & 12 & 176{,}204 \\
\texttt{hodge.pdf}       & Pierre Deligne, \emph{The Hodge Conjecture}       & 5  & 146{,}894 \\
\texttt{riemann.pdf}     & E. Bombieri, \emph{Problems of the Millennium:}   & 11 & 159{,}267 \\
                         & \emph{the Riemann Hypothesis}                     &    &           \\
\texttt{yangmills.pdf}   & Jaffe and Witten, \emph{Quantum Yang--Mills}      & 14 & 195{,}976 \\
                         & \emph{Theory}                                     &    &           \\
\texttt{birchswin.pdf}   & Andrew Wiles, \emph{The Birch and}                & 5  & 142{,}908 \\
                         & \emph{Swinnerton-Dyer Conjecture}                 &    &           \\
\texttt{Wilkins-}        & Riemann 1859, \emph{On the Number of Prime}       & 10 & 119{,}504 \\
\texttt{translation.pdf} & \emph{Numbers less than a Given Quantity},         &    &           \\
                         & translated by David R. Wilkins                    &    &           \\
\texttt{riemann1859.pdf} & Not the 1859 paper. See below.                    & 4  & 6{,}911{,}514 \\
\bottomrule
\end{tabular}
\end{center}

\section{The corpus is complete and was not short a file}

Six Millennium problems remain open and the corpus holds five official statements, so the natural
first reading is that Navier--Stokes is missing. That reading is wrong and this work made it before
checking.

The Clay Mathematics Institute index of Millennium Prize problems, read on 2026-09-11, lists five
problems as unsolved -- Birch and Swinnerton-Dyer, Hodge, P versus NP, Riemann, and Yang--Mills and
the mass gap -- and one as solved, Poincar\'e. Navier--Stokes appears in neither list. The corpus
holds exactly the five the index lists as unsolved. It matches the index and is not short anything.

That statement is a fact about how a website categorized its own contents on one date at one URL,
and it is recorded here as that and not as a fact about the status of a problem. A site can be
mid-edit, deliberately cautious, or behind. The Navier--Stokes problem's own page, read the same
day, carries the status label \emph{Active}. The two readings are not consistent with each other,
both are recorded, and neither is treated here as evidence about the mathematics.

\section{The document that is not what it is named}

\texttt{riemann1859.pdf} is not Riemann's 1859 memoir. It is a photographic scan of Riemann's
handwritten working manuscript.

The evidence is in the file. Four pages carry zero extractable characters between them, measured
page by page. The embedded document title is the scanner operator's own filesystem path,
\texttt{Cod\_\allowbreak{}Ms\_\allowbreak{}Riemann\_\allowbreak{}3\_\allowbreak{}jpg} under a
directory named for the Riemann \emph{Nachlass} and a subdirectory named
\emph{Rechenseiten}, calculation pages. The producer chain is a JPEG passed through PDF24 Creator
and a Quartz PDF context in 2015, which is a scan-to-PDF path and not a typesetting one. All four
pages were rendered and read as images. They are German cursive in ink, heavily struck through with
interlinear corrections, foxed and water-stained, with \emph{Nachlass Riemann} and the year 1859
legible in the header of the first page. The zeta function, the factor $\Pi(s-1)$, the logarithmic
expansion of the Euler product, the function $\xi(t)$ and the logarithmic integral $\mathrm{Li}(x)$
are identifiable in the hand.

So the file is a primary source in the most literal available sense, and it is unusable for text
search or for quotation without somebody who reads nineteenth-century German
\emph{Kurrentschrift}. The content of the published memoir is in the Wilkins translation instead.

\textbf{This matters to the citation registry and not only to a reader.} The registry seeds an
entry with empty bibliographic fields and holds them empty until somebody looks the source up. An
entry seeded from this filename would name Riemann's published 1859 memoir and point at a
manuscript draft, and the empty-field convention would not catch it, because the error is in the
identification and not in the completeness. The distinction has to be written into the entry.

\section{An observation about production}

Five of the six official statements share a producer string of MiKTeX pdfTeX-1.21a and creation
timestamps on 2026-08-04 falling within three minutes of one another. Bombieri's Riemann statement
is older tooling, a dvips and Acrobat Distiller chain dated 2000-08-30, and carries the embedded
title \texttt{Riemann-new.dvi}.

The Navier--Stokes statement, fetched separately from the Clay website for the next chapter, carries
the same producer and a timestamp of 16:46:57, nineteen seconds after the Birch and
Swinnerton-Dyer file. It is part of the same batch.

This is a real observation about how the statements were produced and it says nothing whatever
about their contents. It is recorded because it establishes that the corpus and the separately
fetched statement are the same edition of the same set, which is worth knowing in a registry entry.
It is worth nothing in an argument about mathematics and is not used as one.
